\documentclass[5p]{elsarticle}
\usepackage{amsmath}
\usepackage{booktabs}
\usepackage{graphicx}
\usepackage{xurl}
\journal{Economics Letters}

\begin{document}

\begin{frontmatter}

\title{Entry Regulation and the Allocation of Cross-Border Cannabis Demand}

\author[1]{Mahyar Rezazadeh}
\author[1]{Babak Sahragard}
\address[1]{Department of Economics}

\begin{abstract}
Counties near a still-prohibited neighboring state show a robust cannabis border premium: +48\% revenue per
store, +193\% revenue per capita, +101\% dispensary density (five-state panel, 2014--2026, time-varying border
classification). Capped entry -- coded from actual local-government licensing authority, not state labels --
cuts the per-capita premium's multiplier by 54\% and density by 56\% in relative terms ($p{=}0.009$,
$p{<}0.001$), but leaves the
per-store premium unchanged (3.5\%, $p{=}0.90$). Regulation blocks the quantity margin, not price. A firm-level
panel of 1{,}249 individual New Mexico dispensaries confirms the border premium itself is not a
county-aggregation artifact (+141\%, $p{=}0.001$, significant every year 2022--2026).
\end{abstract}

\begin{keyword}
cannabis \sep cross-border shopping \sep market entry \sep license caps \sep entry regulation
\end{keyword}

\end{frontmatter}

\noindent\textbf{Highlights}
\begin{itemize}
\item Cannabis border counties show higher revenue per store (+48\%), per capita (+193\%), and density
(+101\%) than interior counties, pooling five states with time-varying border exposure.
\item Capped entry -- including states with no state cap but binding local-government control over licensing
-- cuts the border premium's multiplier on revenue per capita (54\%) and dispensary density (56\%) in relative terms.
\item The premium on revenue per store does not differ by entry regime at all.
\item Oregon's own 2022 entry freeze produces the same-signed within-state pattern.
\item Colorado and Montana are reclassified as capped despite no state numerical ceiling, since local
governments control licensing there as effectively as a cap would.
\item A dispensary-level panel for New Mexico (1{,}249 firms, individual addresses and revenue) confirms the
border premium at the level of the individual seller, not just as a county aggregate.
\end{itemize}

\section{Introduction}

Uneven state-level cannabis legalization creates a spatial demand spillover: counties near a still-prohibited
neighboring state face captive cross-border demand that interior counties of the same state do not. Whether
this shows up as higher revenue at existing dispensaries or as more dispensaries serving more of the market is
not a fact about demand -- it is a fact about entry regulation. A capped-entry state closes off the
store-count channel by construction; a free-entry state leaves it open. We test this directly using a
five-state panel and an entry-regime classification built from actual statutes, not labels.

\citet{HaoCowan2020} and \citet{HansenMillerWeber2020} document that cannabis demand crosses legal/nonlegal
state borders; \citet{HollenbeckUetake2021} and \citet{Thomas2019LicenseQuotas} document that cannabis entry
caps distort market structure and rents. No paper connects the two -- asking where captured cross-border
demand goes once supply-side rules determine which margin can respond.

\section{Data and Empirical Strategy}

We use county-year revenue and dispensary-count data for Colorado, Washington, Oregon, Montana, and New
Mexico (2014--2026; each state's own revenue agency, detailed sourcing in the working paper). Border exposure
is a 70-mile indicator, recomputed \emph{by year} rather than fixed to a single snapshot, since New Mexico
(2022), Oregon (2015), and Oregon's own relationship to California/Nevada (2017--2018) each legalize partway
through a neighboring sample state's panel -- this reclassifies 41 of 1{,}314 county-years relative to a
static classification. Entry regime is coded \texttt{Capped} at the state-year level from actual statutes:
Colorado and Montana are capped despite no state numerical ceiling, since local governments bind licensing
directly (48.4\% of Colorado counties ban retail entirely, \citealp{PayanBrownSong2021}; Montana requires
county opt-in from launch); Washington caps by statute; Oregon is free 2015--2021 and capped 2022--2025 (a
genuine within-state switch); New Mexico is free throughout (statutory anti-cap provision). We estimate
$y_{ist} = \alpha + \beta_1 \text{Border}_{ist} + \beta_2 \text{Capped}_{st} + \beta_3 \text{Border}_{ist}
\times \text{Capped}_{st} + \gamma\ln(\text{pop}_{ist}) + \delta_t + \eta_i + \varepsilon_{ist}$ for
log(revenue per store), log(revenue per capita), and log(dispensary density), county-clustered and HC1
standard errors. As a firm-level robustness check, we additionally build a genuine dispensary-level panel for
New Mexico -- the only sample state with historical individual-firm revenue -- from 51 monthly licensee-level
files, identifying each dispensary by its normalized street address and its entry month as the first month
that address appears (1{,}253 dispensaries; 231 left-censored as pre-existing at the panel's April 2022 start;
county assignment via the same ZIP-to-county crosswalk used for the county-level panel, not a new geocoding
pipeline). Because New Mexico's border exposure and entry regime are both time-invariant in-sample, this panel
cannot test the entry-regulation interaction, only whether the border premium itself survives at the firm
level.

\section{Results}

\input{tables/table_main_pooled_el}

Table~\ref{tab:main_pooled} shows a robust border premium on all three outcomes. The interaction with capped
entry is a precisely estimated null on revenue per store (3.5\%, $p{=}0.90$) but a large, significant
suppression on revenue per capita ($-54$ points, $p{=}0.009$) and dispensary density ($-56$ points,
$p{<}0.001$) -- nearly identical magnitudes, consistent with the entire capped-entry suppression running
through the quantity margin, since per-capita revenue is arithmetically close to density times per-store
revenue. State-specific regressions corroborate this at the state level: Colorado's density coefficient is
negative ($-27\%$, $p{<}0.01$) despite a large positive per-store premium ($+81\%$), and New Mexico's density
coefficient is strongly positive ($+109\%$). Oregon's own within-state before/after split around its January
2022 freeze shows the same direction (per-capita premium $+173\%$ free vs.\ $+106\%$ capped) though the formal
interaction is underpowered on three border counties alone ($p{=}0.50$).

A dispensary-level re-estimation on New Mexico's individual sellers (Figure~\ref{fig:nm_dispensary}) recovers a
firm-level border premium of $+141.4\%$ ($p{=}0.001$), significant in every sample year (113--184\%), confirming
the county-level result is not an aggregation artifact -- though New Mexico's own entry regime never varies
in-sample, so this check cannot speak to the capped/free interaction itself. A companion cut isolating
\emph{which level of government} imposes the restriction (\texttt{LocalBarrier}: Colorado and Montana coded 1,
Washington/Oregon/New Mexico coded 0) preserves the density interaction under county clustering ($-41.5\%$,
$p{=}0.037$) but not the per-capita or per-store interactions ($p{=}0.35$, $p{=}0.17$), a caveat on the main
classification detailed in the working paper.

\begin{figure}[htbp]
\centering
\includegraphics[width=0.6\textwidth]{figures/fig13_nm_dispensary_eventstudy.png}
\caption{Border premium on individual-dispensary revenue, by year, New Mexico firm-level panel
($n=1{,}249$ dispensaries).}
\label{fig:nm_dispensary}
\end{figure}

\section{Conclusion}

Cross-border cannabis demand generates a robust revenue premium in border counties. Entry regulation --
classified by actual local-government licensing authority rather than a state's own numerical-cap label --
determines whether that premium flows into market expansion (revenue per capita, dispensary count) or not,
but does not redirect it into the per-store margin the way a simple capped-versus-free story would predict.
The per-store premium is the same size regardless of regime. This is a narrower finding than the paper's
starting hypothesis, and the more defensible one for it.

\bibliographystyle{elsarticle-harv}
\bibliography{references}

\end{document}
